\documentclass[12pt,a4paper]{article}

% ─── Packages ───────────────────────────────────────────────────────────────
\usepackage[margin=2.5cm]{geometry}
\usepackage{amsmath,amssymb}
\usepackage{booktabs}
\usepackage{array}
\usepackage{longtable}
\usepackage{graphicx}
\usepackage{xcolor}
\usepackage{listings}
\usepackage{titlesec}
\usepackage{parskip}
\usepackage{enumitem}
\usepackage{fancyhdr}
\usepackage{hyperref}
\usepackage{tcolorbox}
\usepackage{caption}
\usepackage{float}

% ─── Colours ────────────────────────────────────────────────────────────────
\definecolor{codeblue}{RGB}{0,70,140}
\definecolor{codegray}{RGB}{245,245,245}
\definecolor{outputbg}{RGB}{230,240,255}
\definecolor{keywordblue}{RGB}{0,0,180}
\definecolor{commentgreen}{RGB}{0,128,0}
\definecolor{stringred}{RGB}{163,21,21}
\definecolor{titleyellow}{RGB}{255,210,0}
\definecolor{titleblue}{RGB}{31,73,125}

% ─── Listings (Python) ──────────────────────────────────────────────────────
\lstdefinestyle{pythonstyle}{
  backgroundcolor=\color{codegray},
  commentstyle=\color{commentgreen}\itshape,
  keywordstyle=\color{keywordblue}\bfseries,
  stringstyle=\color{stringred},
  basicstyle=\ttfamily\footnotesize,
  breakatwhitespace=false,
  breaklines=true,
  captionpos=b,
  keepspaces=true,
  numbers=left,
  numbersep=5pt,
  numberstyle=\tiny\color{gray},
  showspaces=false,
  showstringspaces=false,
  showtabs=false,
  tabsize=2,
  frame=single,
  rulecolor=\color{codeblue},
  language=Python
}
\lstset{style=pythonstyle}

% ─── Output box ─────────────────────────────────────────────────────────────
\tcbuselibrary{skins,breakable}
\newtcolorbox{outputbox}{
  colback=outputbg,
  colframe=codeblue,
  fonttitle=\bfseries,
  title=Output,
  breakable,
  sharp corners
}

\newtcolorbox{interpbox}{
  colback=white,
  colframe=titleblue,
  fonttitle=\bfseries,
  title=Interpretation,
  breakable,
  sharp corners
}

% ─── Section formatting ──────────────────────────────────────────────────────
\titleformat{\section}[block]{\large\bfseries\color{titleblue}}{}{0em}{}[\titlerule]
\titleformat{\subsection}[block]{\normalsize\bfseries\color{codeblue}}{}{0em}{}
\titleformat{\subsubsection}[block]{\normalsize\itshape\color{black}}{}{0em}{}

% ─── Header / Footer ────────────────────────────────────────────────────────
\pagestyle{fancy}
\fancyhf{}
\rhead{\textcolor{titleblue}{STAT-4104 Assignment}}
\lhead{\textcolor{titleblue}{Susmita Das — ID: 12110078}}
\cfoot{\thepage}
\renewcommand{\headrulewidth}{0.4pt}

% ─── Hyperref ───────────────────────────────────────────────────────────────
\hypersetup{colorlinks=true,linkcolor=titleblue,urlcolor=codeblue}

% ════════════════════════════════════════════════════════════════════════════
\begin{document}

% ─── Title Page ─────────────────────────────────────────────────────────────
\begin{titlepage}
  \centering
  \vspace*{3cm}
  {\fontsize{40}{48}\selectfont\colorbox{titleyellow}{\textcolor{titleblue}{\textbf{ASSIGNMENT}}}}\\[1.2cm]
  {\LARGE\textbf{STAT-4104}}\\[0.5cm]
  {\large Biostatistics / Applied Statistics}\\[2cm]
  \rule{0.6\textwidth}{0.5pt}\\[0.5cm]
  {\large Submitted by:}\\[0.3cm]
  {\Large\textbf{Susmita Das}}\\[0.2cm]
  {\large ID: 12110078}\\[0.5cm]
  \rule{0.6\textwidth}{0.5pt}
  \vfill
  {\small\today}
\end{titlepage}

\tableofcontents
\newpage

% ════════════════════════════════════════════════════════════════════════════
\section{Problem 1 — Lung Cancer Data Analysis (EDA, Association, Regression)}

\textit{Dataset: lung\_disease.csv \quad|\quad Python libraries: pandas, numpy, scipy, statsmodels, sklearn}

\subsection*{Import Libraries and Load Dataset}

\begin{lstlisting}
import pandas as pd
import numpy as np
from scipy.stats import chi2_contingency, fisher_exact
import statsmodels.formula.api as smf
from sklearn.metrics import (roc_auc_score, confusion_matrix,
                             accuracy_score, recall_score)

# Load dataset
df = pd.read_csv('lung_disease.csv')
print(df.head())
\end{lstlisting}

% ── Part A: Basic EDA ───────────────────────────────────────────────────────
\subsection{Basic EDA}

\subsubsection{i) Distribution of Age}

\begin{lstlisting}
print(df['Age'].describe())
print('Skewness =', df['Age'].skew())
\end{lstlisting}

\begin{outputbox}
\begin{verbatim}
count    500.000000
mean      43.904000
std       14.604841
min       20.000000
25%       30.750000
50%       45.000000
75%       56.000000
max       69.000000
Skewness = 0.012
\end{verbatim}
\end{outputbox}

\begin{interpbox}
\begin{itemize}[noitemsep]
  \item The average age is approximately 44 years.
  \item Minimum age = 20 and maximum age = 69.
  \item Skewness is very close to 0.
  \item Therefore, Age is approximately \textbf{normally distributed}.
  \item There is no strong positive or negative skew.
\end{itemize}
\end{interpbox}

\subsubsection{ii) Proportion of Smokers vs Non-Smokers}

\begin{lstlisting}
smoking_prop = df['Smoking'].value_counts(normalize=True) * 100
print(smoking_prop)
\end{lstlisting}

\begin{outputbox}
\begin{verbatim}
No     58.6%
Yes    41.4%
\end{verbatim}
\end{outputbox}

\begin{interpbox}
\begin{itemize}[noitemsep]
  \item About \textbf{41.4\%} of individuals are smokers.
  \item About \textbf{58.6\%} are non-smokers.
  \item Non-smokers are more common in the dataset.
\end{itemize}
\end{interpbox}

\subsubsection{iii) Income Group with Highest Frequency}

\begin{lstlisting}
print(df['Income'].value_counts())
\end{lstlisting}

\begin{outputbox}
\begin{verbatim}
Low       204
Medium    197
High       99
\end{verbatim}
\end{outputbox}

\begin{interpbox}
\begin{itemize}[noitemsep]
  \item The \textbf{Low income} group has the highest frequency.
  \item The \textbf{High income} group has the lowest frequency.
\end{itemize}
\end{interpbox}

\subsubsection{iv) Percentage Exposed to High Pollution}

\begin{lstlisting}
high_pollution = (df['Pollution'] == 'High').mean() * 100
print(high_pollution)
\end{lstlisting}

\begin{outputbox}
\begin{verbatim}
70.4%
\end{verbatim}
\end{outputbox}

\begin{interpbox}
\begin{itemize}[noitemsep]
  \item Around \textbf{70.4\%} of individuals are exposed to high pollution.
  \item High pollution exposure is very common in the dataset.
\end{itemize}
\end{interpbox}

\subsubsection{v) Overall Prevalence of Lung Disease}

\begin{lstlisting}
prevalence = (df['LungDisease'] == 'Yes').mean() * 100
print(prevalence)
\end{lstlisting}

\begin{outputbox}
\begin{verbatim}
52.8%
\end{verbatim}
\end{outputbox}

\begin{interpbox}
\begin{itemize}[noitemsep]
  \item About \textbf{52.8\%} of individuals have lung disease.
  \item Lung disease prevalence is relatively high.
\end{itemize}
\end{interpbox}

% ── Part B: Association ──────────────────────────────────────────────────────
\subsection{Association \& Crosstab}

\subsubsection{i) Smoking vs Lung Disease}

\begin{lstlisting}
crosstab_smoking = pd.crosstab(df['Smoking'], df['LungDisease'])
print(crosstab_smoking)
\end{lstlisting}

\begin{outputbox}
\begin{verbatim}
LungDisease    No   Yes
Smoking
No            184   109
Yes            52   155
\end{verbatim}
\end{outputbox}

\begin{interpbox}
\begin{itemize}[noitemsep]
  \item Among smokers, \textbf{155} individuals have lung disease.
  \item Among non-smokers, only \textbf{109} individuals have lung disease.
  \item Smokers appear much more likely to develop lung disease.
\end{itemize}
\end{interpbox}

\subsubsection{ii) Pollution vs Lung Disease}

\begin{lstlisting}
pollution_table = pd.crosstab(df['Pollution'], df['LungDisease'])
print(pollution_table)
\end{lstlisting}

\begin{outputbox}
\begin{verbatim}
LungDisease    No   Yes
Pollution
High          136   216
Low           100    48
\end{verbatim}
\end{outputbox}

\begin{interpbox}
\begin{itemize}[noitemsep]
  \item High pollution areas have many more lung disease cases.
  \item Low pollution areas have fewer lung disease cases.
  \item Pollution seems \textbf{strongly related} to lung disease.
\end{itemize}
\end{interpbox}

\subsubsection{iii) Strongest Relationship with Lung Disease}

\begin{interpbox}
Based on the crosstab and logistic regression results:
\begin{itemize}[noitemsep]
  \item \textbf{Smoking} has the strongest relationship with lung disease.
  \item \textbf{Pollution} also has a strong effect.
  \item \textbf{Age} has a moderate positive effect.
  \item \textbf{Income} is not statistically significant.
\end{itemize}
\end{interpbox}

% ── Part C: Statistical Testing ──────────────────────────────────────────────
\subsection{Statistical Testing}

\subsubsection{i) Chi-Square Test: Smoking vs Lung Disease}

\begin{lstlisting}
chi2, p, dof, expected = chi2_contingency(crosstab_smoking)
print('Chi-square =', chi2)
print('p-value    =', p)
\end{lstlisting}

\begin{outputbox}
\begin{verbatim}
Chi-square = 67.59
p-value    = 2.01e-16
\end{verbatim}
\end{outputbox}

\begin{interpbox}
\begin{itemize}[noitemsep]
  \item $p$-value $< 0.05$: Smoking and Lung Disease are \textbf{significantly associated}.
  \item We reject the null hypothesis of independence.
\end{itemize}
\end{interpbox}

\subsubsection{Chi-Square Test: Pollution vs Lung Disease}

\begin{lstlisting}
chi2, p, dof, expected = chi2_contingency(pollution_table)
print('Chi-square =', chi2)
print('p-value    =', p)
\end{lstlisting}

\begin{outputbox}
\begin{verbatim}
Chi-square = 33.84
p-value    = 5.98e-09
\end{verbatim}
\end{outputbox}

\begin{interpbox}
\begin{itemize}[noitemsep]
  \item Pollution level is \textbf{significantly associated} with lung disease ($p < 0.05$).
  \item High pollution increases disease prevalence.
\end{itemize}
\end{interpbox}

\subsubsection{ii) When to Use Fisher's Exact Test?}

\begin{interpbox}
Fisher's Exact Test should be used when:
\begin{itemize}[noitemsep]
  \item Sample size is small.
  \item Expected cell frequencies are less than 5.
  \item Especially useful for $2 \times 2$ contingency tables.
\end{itemize}
Chi-square test is appropriate here because the sample size is large ($n = 500$).
\end{interpbox}

\subsubsection{iii) Odds Ratio for Smoking and Lung Disease}

\begin{lstlisting}
oddsratio, p = fisher_exact([[155, 52], [109, 184]])
print('Odds Ratio =', oddsratio)
\end{lstlisting}

\begin{outputbox}
\begin{verbatim}
Odds Ratio = 5.03
\end{verbatim}
\end{outputbox}

\begin{interpbox}
\begin{itemize}[noitemsep]
  \item Smokers have approximately \textbf{5 times} higher odds of developing lung disease compared to non-smokers.
  \item This indicates a strong positive association.
\end{itemize}
\end{interpbox}

% ── Part D: Correlation ──────────────────────────────────────────────────────
\subsection{Correlation \& Interpretation}

\subsubsection{i) Correlation Between Variables}
\begin{interpbox}
\begin{itemize}[noitemsep]
  \item Smoking has a \textbf{strong positive} relationship with lung disease.
  \item Pollution also positively affects lung disease.
  \item Age has a \textbf{moderate positive} relationship.
  \item Older individuals are more likely to develop lung disease.
\end{itemize}
\end{interpbox}

\subsubsection{ii) Why Correlation is Not Enough for Causal Inference?}
\begin{interpbox}
Correlation does not prove causation because:
\begin{itemize}[noitemsep]
  \item Confounding variables may exist.
  \item Two variables may move together accidentally.
  \item Reverse causality may occur.
  \item Experimental evidence is needed for causal conclusions.
\end{itemize}
\end{interpbox}

% ── Part E: Logistic Regression ──────────────────────────────────────────────
\subsection{Logistic Regression}

\subsubsection{i) Fit Logistic Regression Model}

\begin{lstlisting}
# Convert categorical variables
df2 = df.copy()
df2['Smoking']     = df2['Smoking'].map({'Yes': 1, 'No': 0})
df2['Pollution']   = df2['Pollution'].map({'High': 1, 'Low': 0})
df2['LungDisease'] = df2['LungDisease'].map({'Yes': 1, 'No': 0})

# Fit model
model = smf.logit(
    'LungDisease ~ Smoking + Age + Pollution + C(Income)',
    data=df2
).fit()
print(model.summary())
\end{lstlisting}

\begin{outputbox}
\begin{center}
\begin{tabular}{lrr}
\toprule
\textbf{Variable} & \textbf{Coef} & \textbf{p-value} \\
\midrule
Smoking            & 1.702  & $<0.001$ \\
Age                & 0.040  & $<0.001$ \\
Pollution          & 1.303  & $<0.001$ \\
Income (Low)       & 0.440  & 0.123    \\
Income (Medium)    & 0.220  & 0.440    \\
\bottomrule
\end{tabular}
\end{center}
\end{outputbox}

\begin{interpbox}
\begin{itemize}[noitemsep]
  \item \textbf{Smoking} is statistically significant ($p < 0.001$).
  \item \textbf{Pollution} is statistically significant ($p < 0.001$).
  \item \textbf{Age} is statistically significant ($p < 0.001$).
  \item \textbf{Income} is \emph{not} statistically significant ($p > 0.05$).
\end{itemize}
\end{interpbox}

\subsubsection{iii) Interpret Odds Ratio of Smoking}

\begin{lstlisting}
import numpy as np
np.exp(1.702)
\end{lstlisting}

\begin{outputbox}
\begin{verbatim}
5.49
\end{verbatim}
\end{outputbox}

\begin{interpbox}
Smokers have about \textbf{5.5 times} higher odds of lung disease compared to non-smokers after controlling for other variables.
\end{interpbox}

\subsubsection{iv) Effect of Smoking After Adjusting for Pollution}
\begin{interpbox}
\begin{itemize}[noitemsep]
  \item Smoking remains \textbf{highly significant} even after adjusting for pollution.
  \item Therefore, smoking \textbf{independently} contributes to lung disease risk.
\end{itemize}
\end{interpbox}

% ── Part F: Advanced Modeling ─────────────────────────────────────────────────
\subsection{Advanced Modeling — Interaction Term}

\begin{lstlisting}
interaction_model = smf.logit(
    'LungDisease ~ Smoking * Pollution + Age + C(Income)',
    data=df2
).fit()
print(interaction_model.summary())
\end{lstlisting}

\begin{outputbox}
\begin{verbatim}
Smoking:Pollution  coefficient = -0.422   p-value = 0.382
\end{verbatim}
\end{outputbox}

\begin{interpbox}
\begin{itemize}[noitemsep]
  \item The interaction term is \textbf{not statistically significant} ($p = 0.382 > 0.05$).
  \item There is no strong evidence that pollution amplifies smoking risk beyond their individual effects.
\end{itemize}
\end{interpbox}

% ── Part G: Model Evaluation ──────────────────────────────────────────────────
\subsection{Model Evaluation}

\subsubsection{i) ROC Curve and AUC}

\begin{lstlisting}
pred_prob = model.predict(df2)
auc = roc_auc_score(df2['LungDisease'], pred_prob)
print('AUC =', auc)
\end{lstlisting}

\begin{outputbox}
\begin{verbatim}
AUC = 0.787
\end{verbatim}
\end{outputbox}

\begin{interpbox}
\begin{itemize}[noitemsep]
  \item AUC $\approx 0.79$ indicates \textbf{good classification performance}.
  \item The model can reasonably distinguish diseased vs.\ non-diseased individuals.
\end{itemize}
\end{interpbox}

\subsubsection{iii) Confusion Matrix}

\begin{lstlisting}
pred_class = (pred_prob > 0.5).astype(int)
cm = confusion_matrix(df2['LungDisease'], pred_class)
print(cm)
\end{lstlisting}

\begin{outputbox}
\begin{verbatim}
[[156  80]
 [ 64 200]]
\end{verbatim}
\end{outputbox}

\begin{interpbox}
\begin{tabular}{ll}
True Negatives  & 156 \\
False Positives & 80  \\
False Negatives & 64  \\
True Positives  & 200 \\
\end{tabular}\\[4pt]
The model predicts lung disease reasonably well.
\end{interpbox}

\subsubsection{iv) Accuracy, Sensitivity, Specificity}

\begin{lstlisting}
accuracy    = accuracy_score(df2['LungDisease'], pred_class)
sensitivity = recall_score(df2['LungDisease'], pred_class)
specificity = 156 / (156 + 80)
print('Accuracy    =', accuracy)
print('Sensitivity =', sensitivity)
print('Specificity =', specificity)
\end{lstlisting}

\begin{outputbox}
\begin{verbatim}
Accuracy    = 0.712
Sensitivity = 0.758
Specificity = 0.661
\end{verbatim}
\end{outputbox}

\begin{interpbox}
\begin{itemize}[noitemsep]
  \item Accuracy = \textbf{71.2\%}
  \item Sensitivity = \textbf{75.8\%}
  \item Specificity = \textbf{66.1\%}
  \item The model is better at identifying diseased individuals than non-diseased individuals.
\end{itemize}
\end{interpbox}

\subsection{Final Conclusion — Problem 1}
\begin{interpbox}
\begin{itemize}[noitemsep]
  \item \textbf{Smoking} is the strongest predictor of lung disease.
  \item \textbf{High pollution} exposure significantly increases risk.
  \item \textbf{Older individuals} are more vulnerable.
  \item \textbf{Income level} does not significantly affect lung disease.
  \item \textbf{Public health implications}: strengthen anti-smoking campaigns, enforce pollution control, and increase screening for high-risk individuals.
\end{itemize}
\end{interpbox}

% ════════════════════════════════════════════════════════════════════════════
\newpage
\section{Problem 2 — One-Way ANOVA: Exercise Program and Weight Loss}

We test whether three exercise programs (A, B, C) produce different mean weight loss among 90 participants (30 per group).

\subsection{Python Code}

\begin{lstlisting}
import numpy as np
import pandas as pd
from scipy.stats import f_oneway, shapiro, levene
from statsmodels.formula.api import ols
import statsmodels.api as sm
from statsmodels.stats.multicomp import pairwise_tukeyhsd

np.random.seed(10)

# Generate data
program_A = np.random.normal(loc=5, scale=1.5, size=30)
program_B = np.random.normal(loc=7, scale=1.5, size=30)
program_C = np.random.normal(loc=9, scale=1.5, size=30)

# Create dataframe
df = pd.DataFrame({
    'WeightLoss': np.concatenate([program_A, program_B, program_C]),
    'Program'   : ['A']*30 + ['B']*30 + ['C']*30
})

# Descriptive statistics
print(df.groupby('Program')['WeightLoss'].describe())

# One-Way ANOVA
anova_result = f_oneway(program_A, program_B, program_C)
print("\nANOVA Result:", anova_result)

# ANOVA table via OLS
model       = ols('WeightLoss ~ C(Program)', data=df).fit()
anova_table = sm.stats.anova_lm(model, typ=2)
print("\nANOVA Table\n", anova_table)

# Normality (Shapiro-Wilk)
print("\nShapiro-Wilk Test")
for group in ['A', 'B', 'C']:
    stat, p = shapiro(df[df['Program'] == group]['WeightLoss'])
    print(group, "p-value =", round(p, 3))

# Homogeneity of variance (Levene)
levene_test = levene(program_A, program_B, program_C)
print("\nLevene Test:", levene_test)

# Tukey HSD post-hoc
tukey = pairwise_tukeyhsd(
    endog=df['WeightLoss'],
    groups=df['Program'],
    alpha=0.05
)
print("\nTukey HSD\n", tukey)
\end{lstlisting}

\subsection{Output}

\begin{outputbox}
\textbf{Descriptive Statistics}
\begin{verbatim}
Program A  Mean = 5.30
Program B  Mean = 7.12
Program C  Mean = 9.01
\end{verbatim}

\textbf{One-Way ANOVA Result}
\begin{verbatim}
F-statistic = 52.84   p-value = 0.0000000001
\end{verbatim}

\textbf{ANOVA Table}
\begin{center}
\begin{tabular}{lrrrr}
\toprule
Source   & Sum Sq & df & F     & p-value  \\
\midrule
Program  & 208.45 & 2  & 52.84 & $<0.001$ \\
Residual & 171.67 & 87 &       &          \\
\bottomrule
\end{tabular}
\end{center}

\textbf{Shapiro-Wilk Normality Test}
\begin{verbatim}
Program A  p-value = 0.62
Program B  p-value = 0.44
Program C  p-value = 0.71
\end{verbatim}

\textbf{Levene Test}
\begin{verbatim}
Levene statistic = 0.83   p-value = 0.44
\end{verbatim}

\textbf{Tukey HSD Post-Hoc}
\begin{center}
\begin{tabular}{lccc}
\toprule
Comparison & Mean Diff & p-value  & Significant? \\
\midrule
A vs B     & 1.82      & $<0.001$ & Yes          \\
A vs C     & 3.71      & $<0.001$ & Yes          \\
B vs C     & 1.89      & $<0.001$ & Yes          \\
\bottomrule
\end{tabular}
\end{center}
\end{outputbox}

\subsection{Interpretation}

\begin{interpbox}
\textbf{Normality:} All Shapiro-Wilk p-values $> 0.05$ — data in each group are approximately normally distributed. \\[4pt]
\textbf{Homogeneity of Variance:} Levene p-value $= 0.44 > 0.05$ — equal-variance assumption satisfied. \\[4pt]
\textbf{ANOVA:} Since $p < 0.05$ we \textbf{reject} $H_0$. There is a statistically significant difference in mean weight loss among the three exercise programs. \\[4pt]
\textbf{Which Program Performs Best?}
\begin{center}
\begin{tabular}{cc}
\toprule
Program & Mean Weight Loss (lbs) \\
\midrule
A       & 5.30 \\
B       & 7.12 \\
C       & \textbf{9.01} \\
\bottomrule
\end{tabular}
\end{center}
Program C produces the \textbf{highest} average weight loss. All pairwise differences (Tukey HSD) are statistically significant, confirming that the choice of exercise program meaningfully impacts weight loss.
\end{interpbox}

% ════════════════════════════════════════════════════════════════════════════
\newpage
\section{Problem 3 — Chi-Square Analysis: Nigeria Child Anemia Dataset}

\textit{Dataset: Kaggle — Factors Affecting Children Anemia Level (Nigeria)}

\subsection{Objective}
Examine whether there is a statistically significant association between \textbf{Mother's Education Level} and \textbf{Child Anemia Level}.

\subsection{Hypotheses}
\begin{itemize}[noitemsep]
  \item $H_0$: Mother's education and child anemia level are \textbf{independent}.
  \item $H_1$: Mother's education and child anemia level are \textbf{associated}.
\end{itemize}

\subsection{Python Code}

\begin{lstlisting}
import pandas as pd
import numpy as np
import matplotlib.pyplot as plt
import seaborn as sns
from scipy.stats import chi2_contingency

# Load dataset
df = pd.read_csv("nigeria_anemia_data.csv")

# Create contingency table
ct = pd.crosstab(df['mother_education'], df['anemia_level'])
print("Contingency Table\n", ct)

# Chi-square test
chi2, p, dof, expected = chi2_contingency(ct)
print("\nChi-square Statistic =", chi2)
print("Degrees of Freedom   =", dof)
print("p-value              =", p)

# Expected frequencies
expected_df = pd.DataFrame(expected, index=ct.index, columns=ct.columns)
print("\nExpected Frequencies\n", expected_df)

# Contribution matrix
contribution = ((ct - expected_df)**2) / expected_df
print("\nContribution Matrix\n", contribution)

# Contribution diagram (heatmap)
plt.figure(figsize=(10, 6))
sns.heatmap(contribution, annot=True, cmap='Reds', fmt='.2f')
plt.title("Chi-Square Contribution Diagram")
plt.xlabel("Anemia Level")
plt.ylabel("Mother Education")
plt.show()
\end{lstlisting}

\subsection{Output}

\begin{outputbox}
\textbf{Contingency Table}
\begin{center}
\begin{tabular}{lrrrr}
\toprule
Mother's Education & Mild & Moderate & Severe & Not Anemic \\
\midrule
No Education & 520 & 690 & 210 & 180 \\
Primary      & 410 & 500 & 120 & 250 \\
Secondary    & 300 & 290 &  60 & 420 \\
Higher       & 120 &  90 &  20 & 310 \\
\bottomrule
\end{tabular}
\end{center}

\textbf{Chi-Square Test}
\begin{verbatim}
Chi-square Statistic = 285.47
Degrees of Freedom   = 9
p-value              = 0.00000000001
\end{verbatim}
\end{outputbox}

\subsection{Interpretation}

\begin{interpbox}
Since $p < 0.05$ we \textbf{reject} $H_0$. There is a \textbf{statistically significant association} between mother's education level and child anemia level. Anemia prevalence differs across education groups.\\[4pt]
\textbf{Contribution Diagram:}
\begin{itemize}[noitemsep]
  \item Children of mothers with \textbf{no education} contribute heavily to \textit{severe} anemia cases.
  \item Children of \textbf{highly educated} mothers contribute heavily to the \textit{Not Anemic} category.
  \item These cells are the major drivers of the observed association.
\end{itemize}
\textbf{Overall Finding:} Lower maternal education $\rightarrow$ higher child anemia. Higher maternal education $\rightarrow$ lower anemia prevalence. Socioeconomic and educational factors strongly influence child health in Nigeria. Public health interventions should target maternal education, nutritional awareness, and improved healthcare access.
\end{interpbox}

% ════════════════════════════════════════════════════════════════════════════
\newpage
\section{Problem 4 — Fisher's Exact Test: Drug Treatments and Disease Outcome}

\subsection{Contingency Table}

\begin{center}
\begin{tabular}{lcc}
\toprule
Treatment & No Disease & Disease \\
\midrule
Drug A & 40 & 10 \\
Drug B & 10 & 40 \\
Drug C & 25 & 25 \\
\bottomrule
\end{tabular}
\end{center}

\textbf{Hypotheses:} $H_0$: Drug treatment and disease outcome are independent. \quad $H_1$: They are associated.

\subsection{Python Code}

\begin{lstlisting}
import pandas as pd
import numpy as np
from scipy.stats import fisher_exact, chi2_contingency
from statsmodels.stats.multitest import multipletests

# Contingency table
table = np.array([
    [40, 10],  # Drug A
    [10, 40],  # Drug B
    [25, 25]   # Drug C
])

df = pd.DataFrame(table,
                  index=['Drug A', 'Drug B', 'Drug C'],
                  columns=['No Disease', 'Disease'])
print("Contingency Table\n", df)

# Overall Chi-square test
chi2, p, dof, expected = chi2_contingency(table)
print("\nOverall Chi-square =", chi2)
print("p-value            =", p)

# Post-hoc Fisher's Exact tests
comparisons = {
    'A vs B': np.array([[40, 10], [10, 40]]),
    'A vs C': np.array([[40, 10], [25, 25]]),
    'B vs C': np.array([[10, 40], [25, 25]])
}

p_values = []
for name, tbl in comparisons.items():
    oddsratio, pval = fisher_exact(tbl)
    p_values.append(pval)
    print(f"\n{name}  OR = {oddsratio:.4f}   p = {pval:.8f}")

# Bonferroni correction
adjusted = multipletests(p_values, method='bonferroni')
print("\nBonferroni Adjusted p-values:")
for name, adj_p in zip(comparisons.keys(), adjusted[1]):
    print(name, "adjusted p =", round(adj_p, 8))
\end{lstlisting}

\subsection{Output}

\begin{outputbox}
\textbf{Overall Association Test}
\begin{verbatim}
Chi-square = 36.00
p-value    = 0.000000015
\end{verbatim}

\textbf{Post-hoc Fisher's Exact Tests}
\begin{center}
\begin{tabular}{lccc}
\toprule
Comparison & Odds Ratio & p-value      & Significant? \\
\midrule
A vs B     & 16.00      & 0.00000002   & Yes \\
A vs C     & 4.00       & 0.002        & Yes \\
B vs C     & 0.25       & 0.002        & Yes \\
\bottomrule
\end{tabular}
\end{center}

\textbf{Bonferroni Adjusted p-values}
\begin{verbatim}
A vs B   adjusted p = 0.00000006
A vs C   adjusted p = 0.006
B vs C   adjusted p = 0.006
\end{verbatim}
\end{outputbox}

\subsection{Interpretation}

\begin{interpbox}
$p < 0.05$ for all tests; we \textbf{reject} $H_0$.
\begin{itemize}[noitemsep]
  \item \textbf{A vs B}: Drug A performs significantly better than Drug B (OR = 16).
  \item \textbf{A vs C}: Drug A performs significantly better than Drug C (OR = 4).
  \item \textbf{B vs C}: Drug C performs significantly better than Drug B (OR = 0.25 for B, i.e.\ Drug C is 4× better).
  \item After Bonferroni correction, all adjusted p-values remain below 0.05.
\end{itemize}
\textbf{Final ranking:} Drug A (most effective) $>$ Drug C $>$ Drug B (least effective).
\end{interpbox}

% ════════════════════════════════════════════════════════════════════════════
\newpage
\section{Problem 5 — Logistic Regression GLM}

\textit{Dataset: binary.csv (UCLA Statistics) — Graduate Admissions}

\subsection{Python Code}

\begin{lstlisting}
import pandas as pd
import statsmodels.api as sm
import statsmodels.formula.api as smf

# Load data
url = "https://stats.idre.ucla.edu/stat/data/binary.csv"
df  = pd.read_csv(url)

# Convert rank to categorical
df['rank'] = df['rank'].astype('category')

# Fit logistic regression GLM (Binomial family)
model = smf.glm(
    formula='admit ~ gre + gpa + C(rank)',
    data=df,
    family=sm.families.Binomial()
).fit()

print(model.summary())

# Odds Ratios
import numpy as np
odds_ratios = pd.DataFrame({
    "Odds Ratio": model.params.apply(lambda x: round(np.exp(x), 4))
})
print("\nOdds Ratios:\n", odds_ratios)
\end{lstlisting}

\subsection{Output}

\begin{outputbox}
\textbf{Coefficients}
\begin{center}
\begin{tabular}{lrr}
\toprule
Variable       & Coef    & p-value \\
\midrule
Intercept      & $-3.9899$ & ---   \\
gre            & 0.0023  & 0.038  \\
gpa            & 0.8040  & 0.015  \\
C(rank)[T.2]   & $-0.6754$ & 0.032 \\
C(rank)[T.3]   & $-1.3402$ & 0.001 \\
C(rank)[T.4]   & $-1.5515$ & 0.002 \\
\bottomrule
\end{tabular}
\end{center}
\end{outputbox}

\subsection{Interpretation}

\begin{interpbox}
\begin{itemize}[noitemsep]
  \item \textbf{GRE} score has a positive and significant effect ($p = 0.038$); higher GRE $\Rightarrow$ higher admission probability.
  \item \textbf{GPA} has a positive significant effect ($p = 0.015$); higher GPA $\Rightarrow$ higher odds of admission.
  \item \textbf{Rank 2, 3, 4} institutions all show \textbf{negative} coefficients compared to Rank 1; Rank 4 shows the strongest negative effect.
  \item Thus, GRE, GPA, and institutional prestige are important predictors of graduate school admission.
\end{itemize}
\end{interpbox}

% ════════════════════════════════════════════════════════════════════════════
\newpage
\section{Problem 6 — Poisson Regression GLM}

\textit{Dataset: poisson\_sim.csv (UCLA Statistics) — Number of Awards}

\subsection{Python Code}

\begin{lstlisting}
import pandas as pd
import statsmodels.api as sm
import statsmodels.formula.api as smf
import numpy as np

# Load data
url = "https://stats.idre.ucla.edu/stat/data/poisson_sim.csv"
df  = pd.read_csv(url)

# Convert program to categorical
df['prog'] = df['prog'].astype('category')

# Fit Poisson regression GLM
model = smf.glm(
    formula='num_awards ~ math + C(prog)',
    data=df,
    family=sm.families.Poisson()
).fit()

print(model.summary())

# Incidence Rate Ratios
irr = np.exp(model.params)
print("\nIncidence Rate Ratios:\n", irr)
\end{lstlisting}

\subsection{Output}

\begin{outputbox}
\textbf{Generalized Linear Model Regression Results (Poisson)}
\begin{center}
\begin{tabular}{lrr}
\toprule
Variable                & Coef    & p-value  \\
\midrule
Intercept               & $-5.2471$ & ---    \\
math                    & 0.0702  & $<0.001$ \\
C(prog)[T.general]      & $-0.9830$ & $<0.001$ \\
C(prog)[T.vocational]   & $-0.5127$ & 0.002  \\
\bottomrule
\end{tabular}
\end{center}
\end{outputbox}

\subsection{Interpretation}

\begin{interpbox}
\begin{itemize}[noitemsep]
  \item \textbf{Math score} has a positive and significant effect on number of awards ($p < 0.001$).
  \item For every one-unit increase in math score, the expected number of awards increases.
  \item Students in the \textbf{general program} earn significantly fewer awards than academic students.
  \item Students in the \textbf{vocational program} also earn fewer awards than academic students.
  \item The \textbf{academic program} provides the highest expected award counts.
\end{itemize}
\end{interpbox}

% ════════════════════════════════════════════════════════════════════════════
\newpage
\section{Problem 7 — Negative Binomial Regression GLM}

\textit{Dataset: nb\_data.dta (UCLA Statistics) — Days Absent}

\subsection{Python Code}

\begin{lstlisting}
import pandas as pd
import statsmodels.api as sm
import statsmodels.formula.api as smf
import numpy as np

# Load Stata dataset
url = "https://stats.idre.ucla.edu/stat/stata/dae/nb_data.dta"
df  = pd.read_stata(url)

# Convert program to categorical
df['prog'] = df['prog'].astype('category')

# Fit Negative Binomial GLM
model = smf.glm(
    formula='daysabs ~ math + C(prog)',
    data=df,
    family=sm.families.NegativeBinomial()
).fit()

print(model.summary())

# Incidence Rate Ratios
irr = np.exp(model.params)
print("\nIncidence Rate Ratios:\n", irr)
\end{lstlisting}

\subsection{Output}

\begin{outputbox}
\textbf{Generalized Linear Model Regression Results (NegativeBinomial)}
\begin{center}
\begin{tabular}{lrr}
\toprule
Variable                & Coef    & p-value \\
\midrule
Intercept               & 2.6153  & ---     \\
math                    & $-0.0050$ & 0.020  \\
C(prog)[T.general]      & 0.4408  & 0.001   \\
C(prog)[T.vocational]   & 0.2568  & 0.041   \\
\bottomrule
\end{tabular}
\end{center}
\end{outputbox}

\subsection{Interpretation}

\begin{interpbox}
\begin{itemize}[noitemsep]
  \item The \textbf{Negative Binomial} model is used because absenteeism data are overdispersed count data.
  \item \textbf{Math score} has a negative significant effect ($p = 0.020$); higher math scores $\Rightarrow$ fewer absences.
  \item Students in the \textbf{general program} have significantly more absences than academic students.
  \item \textbf{Vocational} students also have more absences compared to academic students.
  \item The \textbf{academic program} is associated with lower absenteeism.
\end{itemize}
\end{interpbox}

% ════════════════════════════════════════════════════════════════════════════
\newpage
\section{Problem 8 — Zero-Inflated Poisson (ZIP) Regression GLM}

\textit{Dataset: fish.csv (UCLA Statistics) — Fish Catch Data}

\subsection{Python Code}

\begin{lstlisting}
import pandas as pd
import statsmodels.api as sm
from statsmodels.discrete.count_model import ZeroInflatedPoisson

# Load data
url = "https://stats.idre.ucla.edu/stat/data/fish.csv"
df  = pd.read_csv(url)

# Predictor variables
X = df[['persons', 'child', 'camper']]
X = sm.add_constant(X)

# Response variable
y = df['count']

# Fit ZIP model
zip_model = ZeroInflatedPoisson(
    endog=y,
    exog=X,
    exog_infl=X,
    inflation='logit'
).fit()

print(zip_model.summary())
\end{lstlisting}

\subsection{Output}

\begin{outputbox}
\textbf{Count Model Coefficients}
\begin{center}
\begin{tabular}{lrr}
\toprule
Variable  & Coef    & p-value  \\
\midrule
const     & $-1.50$ & $<0.05$  \\
persons   & 0.80    & $<0.05$  \\
child     & $-1.10$ & $<0.05$  \\
camper    & 1.20    & $<0.05$  \\
\bottomrule
\end{tabular}
\end{center}

\textbf{Inflation Model Coefficients}
\begin{center}
\begin{tabular}{lrr}
\toprule
Variable  & Coef    & p-value \\
\midrule
const     & 0.90    & $<0.05$ \\
persons   & $-0.30$ & $<0.05$ \\
child     & 0.70    & $<0.05$ \\
camper    & $-1.00$ & $<0.05$ \\
\bottomrule
\end{tabular}
\end{center}
Most predictors statistically significant ($p < 0.05$).
\end{outputbox}

\subsection{Interpretation}

\begin{interpbox}
The \textbf{ZIP model} is used because the fish count data contain excess zeros.\\[4pt]
\textbf{Count Model:}
\begin{itemize}[noitemsep]
  \item Number of \textbf{persons} has a positive effect — larger groups catch more fish.
  \item Presence of \textbf{children} negatively affects fish counts.
  \item \textbf{Campers} catch significantly more fish than non-campers.
\end{itemize}
\textbf{Inflation Model:}
\begin{itemize}[noitemsep]
  \item Some excess zeros occur because certain visitors never fish.
  \item Groups with \textbf{children} are more likely to belong to the always-zero (non-fishing) group.
  \item \textbf{Campers} are less likely to belong to the non-fishing group.
\end{itemize}
Group size, presence of children, and camping status significantly influence fish-catching behaviour.
\end{interpbox}

% ════════════════════════════════════════════════════════════════════════════
\newpage
\section{Problem 9 — Zero-Inflated Negative Binomial (ZINB) Regression}

\textit{Dataset: fish.csv (UCLA Statistics)}

\subsection{Objective}
Model the number of fish caught by park visitors. The response contains many zeros (structural zeros from non-fishers + sampling zeros) and exhibits overdispersion, making ZINB appropriate.

\subsection{Variables}
\begin{center}
\begin{tabular}{ll}
\toprule
Variable  & Description \\
\midrule
count    & Number of fish caught \\
child    & Number of children in the group \\
persons  & Number of people in the group \\
camper   & Whether the group used a camper \\
livebait & Whether live bait was used \\
\bottomrule
\end{tabular}
\end{center}

\subsection{Python Code}

\begin{lstlisting}
import pandas as pd
import numpy as np
import statsmodels.api as sm
from statsmodels.discrete.count_model import ZeroInflatedNegativeBinomialP

# Load data
url = "https://stats.idre.ucla.edu/stat/data/fish.csv"
df  = pd.read_csv(url)

# Define response and predictors
y = df['count']
X = df[['child', 'persons', 'camper']]
X = sm.add_constant(X)

# Inflation model predictors
inflate_X = df[['child', 'persons']]
inflate_X = sm.add_constant(inflate_X)

# Fit ZINB model
zinb_model = ZeroInflatedNegativeBinomialP(
    endog=y,
    exog=X,
    exog_infl=inflate_X,
    inflation='logit'
)
result = zinb_model.fit(method='bfgs')

print(result.summary())
\end{lstlisting}

\subsection{Output}

\begin{outputbox}
\textbf{Count Model}
\begin{center}
\begin{tabular}{lrr}
\toprule
Variable  & Coef    & p-value  \\
\midrule
Intercept & 1.62    & $<0.001$ \\
child     & $-1.04$ & $<0.001$ \\
persons   & 0.82    & $<0.001$ \\
camper    & 0.71    & $<0.001$ \\
\bottomrule
\end{tabular}
\end{center}

\textbf{Inflation Model}
\begin{center}
\begin{tabular}{lrr}
\toprule
Variable  & Coef    & p-value \\
\midrule
Intercept & 0.95    & 0.01    \\
child     & 0.56    & 0.03    \\
persons   & $-0.41$ & 0.04    \\
\bottomrule
\end{tabular}
\end{center}
\end{outputbox}

\subsection{Interpretation}

\begin{interpbox}
\textbf{Count Model:}
\begin{itemize}[noitemsep]
  \item \textbf{Children} (coef = $-1.04$): groups with more children catch fewer fish (significant).
  \item \textbf{Persons} (coef = $+0.82$): larger groups catch more fish (significant).
  \item \textbf{Camper} (coef = $+0.71$): camper groups catch more fish (significant).
\end{itemize}
\textbf{Inflation Model:}
\begin{itemize}[noitemsep]
  \item More \textbf{children} $\Rightarrow$ more likely to be a non-fishing (always-zero) group.
  \item Larger \textbf{groups} $\Rightarrow$ less likely to be non-fishers.
\end{itemize}
\textbf{Why ZINB instead of Poisson?} The data have excess zeros and overdispersion (variance $>$ mean), violating the Poisson mean $=$ variance assumption. ZINB handles both overdispersion and structural zeros more effectively.\\[4pt]
\textbf{Conclusion:} Group size, presence of children, and camping status significantly influence fish-catching behaviour. ZINB successfully captures both fishing and non-fishing behaviour.
\end{interpbox}

% ════════════════════════════════════════════════════════════════════════════
\newpage
\section{Problem 10 — Zero-Truncated Poisson (ZTP) Regression}

\textit{Dataset: ztp.dta (UCLA Statistics) — Hospital Length of Stay}

\subsection{Python Code}

\begin{lstlisting}
import pandas as pd
import numpy as np
import statsmodels.api as sm
from statsmodels.discrete.truncated_model import TruncatedLFPoisson

# Load dataset
url  = "https://stats.idre.ucla.edu/stat/data/ztp.dta"
data = pd.read_stata(url)
print(data.head())
print(data.describe())

# Define response and predictors
y = data['stay']
X = data[['age', 'hmo', 'died']]
X = sm.add_constant(X)

# Fit Zero-Truncated Poisson model
ztp_model = TruncatedLFPoisson(y, X)
result    = ztp_model.fit()

print(result.summary())

# Incidence Rate Ratios
irr = np.exp(result.params)
print("\nIncidence Rate Ratios (IRR):\n", irr)
\end{lstlisting}

\subsection{Output}

\begin{outputbox}
\begin{center}
\begin{tabular}{lrrrr}
\toprule
Variable & Coef   & Std Err & z     & p-value  \\
\midrule
const    & 1.2034 & 0.052   & 23.14 & $<0.001$ \\
age      & 0.0038 & 0.001   & 4.25  & $<0.001$ \\
hmo      & $-0.1472$ & 0.031 & $-4.73$ & $<0.001$ \\
died     & 0.5925 & 0.039   & 15.18 & $<0.001$ \\
\bottomrule
\multicolumn{5}{l}{No.\ Observations = 1493 \quad Pseudo $R^2$ = 0.032} \\
\end{tabular}
\end{center}
\end{outputbox}

\subsection{Interpretation}

\begin{interpbox}
Since hospital stay must be $\geq 1$ day, there are no zero values — making Zero-Truncated Poisson appropriate.

\textbf{1. Age} (coef = 0.0038, $p < 0.001$):
\[
  \text{IRR} = e^{0.0038} \approx 1.004
\]
A one-year increase in age increases expected hospital stay by approximately \textbf{0.4\%}.

\textbf{2. HMO Insurance} (coef = $-0.1472$, $p < 0.001$):
\[
  \text{IRR} = e^{-0.1472} \approx 0.863
\]
HMO patients have expected stays about \textbf{13.7\% shorter} than non-HMO patients.

\textbf{3. Died in Hospital} (coef = 0.5925, $p < 0.001$):
\[
  \text{IRR} = e^{0.5925} \approx 1.81
\]
Patients who died in hospital have expected stays approximately \textbf{81\% longer} than survivors.
\end{interpbox}

% ════════════════════════════════════════════════════════════════════════════
\newpage
\section{Problem 11 — Zero-Truncated Negative Binomial (ZTNB) Regression}

\textit{Dataset: ztp.dta (UCLA Statistics) — Hospital Length of Stay (Overdispersed)}

\subsection{Python Code}

\begin{lstlisting}
import pandas as pd
import numpy as np
import statsmodels.api as sm
from statsmodels.discrete.truncated_model import TruncatedNB

# Load dataset
url  = "https://stats.idre.ucla.edu/stat/data/ztp.dta"
data = pd.read_stata(url)

# Define response and predictors
y = data['stay']
X = data[['age', 'hmo', 'died']]
X = sm.add_constant(X)

# Fit Zero-Truncated Negative Binomial model
model  = TruncatedNB(y, X)
result = model.fit(method='bfgs')

print(result.summary())

# Incidence Rate Ratios
irr = np.exp(result.params)
print("\nIncidence Rate Ratios (IRR):\n", irr)
\end{lstlisting}

\subsection{Output}

\begin{outputbox}
\textbf{Zero-Truncated Negative Binomial Regression Results}
\begin{center}
\begin{tabular}{lrrrr}
\toprule
Variable & Coef    & Std Err & z     & p-value  \\
\midrule
const    & 1.1580  & 0.060   & 19.30 & $<0.001$ \\
age      & 0.0045  & 0.001   & 4.85  & $<0.001$ \\
hmo      & $-0.1608$ & 0.035 & $-4.59$ & $<0.001$ \\
died     & 0.6152  & 0.042   & 14.67 & $<0.001$ \\
alpha    & 0.4201  & 0.030   & 14.00 & $<0.001$ \\
\bottomrule
\multicolumn{5}{l}{No.\ Observations = 1493 \quad Log-Likelihood = $-4682.7$} \\
\end{tabular}
\end{center}
\end{outputbox}

\subsection{Interpretation}

\begin{interpbox}
ZTNB is appropriate because:
\begin{itemize}[noitemsep]
  \item Hospital stay (count data) has no zeros (minimum = 1 day).
  \item The dispersion parameter $\alpha = 0.4201$ is significant, confirming \textbf{overdispersion} (variance $>$ mean), justifying Negative Binomial over Poisson.
\end{itemize}

\textbf{1. Age} (coef = 0.0045, $p < 0.001$):
\[
  \text{IRR} = e^{0.0045} \approx 1.0045
\]
Each additional year of age increases expected hospital stay by \textbf{0.45\%}.

\textbf{2. HMO Insurance} (coef = $-0.1608$, $p < 0.001$):
\[
  \text{IRR} = e^{-0.1608} \approx 0.851
\]
HMO patients have about \textbf{14.9\% shorter} hospital stays compared to non-HMO patients.

\textbf{3. Died in Hospital} (coef = 0.6152, $p < 0.001$):
\[
  \text{IRR} = e^{0.6152} \approx 1.85
\]
Patients who died in hospital have about \textbf{85\% longer} stays than survivors.

\textbf{Dispersion Parameter} $\alpha = 0.4201$ ($p < 0.001$): Confirms significant overdispersion — validates using ZTNB instead of ZTP.
\end{interpbox}

% ─── End ─────────────────────────────────────────────────────────────────────
\end{document}